%% =====================================================================
%% Paper 34, §III.21 — Multipole expansion / Gaunt termination
%% Sprint 1 dictionary-completion draft, Track 2
%%
%% Three-axis row (for §IV Table~\ref{tab:projection_axes}):
%%   Projection:               Multipole expansion / Gaunt termination
%%   Variables added:          $L$ (multipole order, integer; ranges
%%                             $0,1,\dots,L_\text{max} = 2 l_\text{max}$)
%%   Dimension:                preserves (Coulomb kernel re-expressed,
%%                             each multipole term carries the same
%%                             $[E]$ dimension as the parent Coulomb
%%                             matrix element)
%%   Transcendental signature: none beyond Layer~1; Gaunt $3j$
%%                             coefficients are exact rationals (live
%%                             in $\mathbb{Q}[\sqrt{2k{+}1}]_k$ via
%%                             \S~\ref{sec:proj_3j}); $R^L$ radial
%%                             integrals are algebraic in the
%%                             $Z_\text{eff}$-rational basis
%%   Layer transition:         Layer~1 $\to$ Layer~1 (angular
%%                             reorganization of the Coulomb kernel,
%%                             prior to any focal-length projection)
%%
%% Cross-references (live in body):
%%   - §III.8 Wigner $3j$ (sec:proj_3j) — the angular content this
%%     projection ultimately calls at every vertex
%%   - Paper~15 §V — split-region Legendre expansion for V_ne at
%%     Level~4 (where the termination was first stated for the
%%     molecule-frame hyperspherical solver)
%%   - Paper~19 §III.C — analytical evaluation via incomplete-gamma
%%     for cross-center V_ne in the balanced coupled architecture
%%     (Sprint CD, April 2026; termination verified explicitly at
%%     $L_\text{max} = 2 l_\text{max}$, 33 terms at $n_\text{max}=2$,
%%     168 at $n_\text{max}=3$, machine precision)
%%   - §III.18 (sec:proj_magnetization_density) — sibling angular-only
%%     projection in the Zemach radial-integrand sense
%%
%% Structural notes:
%%   - The termination is EXACT, not asymptotic — by Wigner-$3j$
%%     triangle inequality, $G(l_1, L, l_2) = 0$ for $L > l_1 + l_2$.
%%     No truncation parameter, no UV cutoff, no convergence series.
%%   - Workhorse projection: used in every multi-electron matrix
%%     element of $V_{ee}$ and every cross-center $V_{ne}$ matrix
%%     element across the codebase.
%%   - Sibling to §III.8 (Wigner $3j$): the multipole expansion is
%%     the coordinate-side projection that decomposes $1/|r_1 - r_2|$
%%     into angular content (Gaunt integrals at each vertex, which
%%     are the §III.8 calls) and radial content ($R^L = \int
%%     r_<^L / r_>^{L+1}$ radial matrix elements). The two projections
%%     compose: multipole expansion calls Wigner~$3j$ at every L.
%% =====================================================================

\subsection{Multipole expansion / Gaunt termination}
\label{sec:proj_multipole_gaunt}
\textbf{Source:} bare two-body Coulomb kernel $1/|\vec{r}_1 - \vec{r}_2|$
(for $V_{ee}$) or $1/|\vec{r} - \vec{R}_B|$ (for cross-center $V_{ne}$)
acting on a product of angular-momentum eigenstates on the Fock-projected
$S^3$ graph.\\
\textbf{Target:} the same kernel re-expressed as a finite sum over
multipoles
\begin{equation*}
  \frac{1}{|\vec{r}_1 - \vec{r}_2|}
  \;=\; \sum_{L=0}^{L_\text{max}}
        \frac{r_<^L}{r_>^{L+1}}\, P_L(\cos\theta_{12}),
\end{equation*}
with the angular content $P_L(\cos\theta_{12})$ contracted against
spherical-harmonic vertices via Gaunt integrals
$G(l_1, L, l_2) = 2 \begin{pmatrix} l_1 & L & l_2 \\ 0 & 0 & 0
\end{pmatrix}^{\!2}$, i.e.\ Wigner-$3j$ angular coupling
(\S~\ref{sec:proj_3j}) at each vertex.\\
\textbf{Variables introduced:} the multipole order $L$ (integer, runs
$0, 1, \dots, L_\text{max} = 2 l_\text{max}$).  No physical parameter
is introduced; $L$ is an internal summation index whose upper bound is
fixed by the angular content of the basis.\\
\textbf{Dimension introduced:} dimension-preserving.  Each multipole
term $r_<^L / r_>^{L+1}$ carries the same $[E]$ scaling as the parent
Coulomb matrix element; the decomposition reorganizes the angular
content of an existing operator without injecting a new $[L]$, $[E]$,
$[M]$, or $[T]$ scale.\\
\textbf{Transcendental signature:} ring-preserving over the Layer~1
algebraic ring.  The Gaunt $3j$ coefficients live in
$\mathbb{Q}[\sqrt{2k{+}1}]_k$ (the Wigner-$3j$ ring of
\S~\ref{sec:proj_3j}); the radial $R^L$ matrix elements are algebraic
in the hydrogenic / $Z_\text{eff}$-rational basis (Slater-type
integrals evaluate to rationals at integer screening and to
incomplete-gamma combinations at fractional screening, per
\S~\ref{sec:proj_3j} sibling derivations and Paper~19 \S~III.C).  No
transcendental is introduced at any step.\\
\textbf{Used in:} essentially every multi-electron matrix element in
the codebase.  Specific instances: Level~4 mol-frame hyperspherical
$V_{ee}$ and $V_{ne}$ coupling (Paper~15 \S~V, split-region Legendre
expansion); cross-center $V_{ne}$ in the balanced coupled architecture
(Paper~19 \S~III.C; \texttt{geovac/shibuya\_wulfman.py}); single-center
Slater $F^k$ and $G^k$ integrals at every level (Paper~7 \S~VI.B,
Paper~12); the Roothaan cross-register $V_{eN}$ kernel at multi-focal
$\lambda_a \ne \lambda_b$ (\texttt{geovac/cross\_register\_vne.py});
the multipole structure underlying \S~\ref{sec:proj_tensor_multipole}
(nuclear rank-$L$ couplings).\\
\textbf{Empirical anchor:} the multipole expansion underwrites the
angular sparsity of the Pauli-string Hamiltonian across the catalogue
(Paper~22 angular sparsity theorem; ERI density $1.44\%$ at
$l_\text{max} = 3$ depends on the exact $L \leq l_1 + l_2$
truncation).  Direct numerical verification in Sprint CD (April 2026,
Paper~19 \S~III.C, Track CD v2.0.39): cross-center $V_{ne}$ for LiH at
$n_\text{max} = 2$ contributes exactly 33 nonzero one-body matrix
elements (multipole sum runs $L = 0, 1, 2$, terminates at $L_\text{max}
= 2$); at $n_\text{max} = 3$, 168 nonzero elements ($L_\text{max} =
4$).  Machine-precision evaluation via the regularized incomplete
gamma function (\texttt{scipy.special.gammainc} / \texttt{gammaincc})
replaces grid-based trapezoid quadrature with zero truncation error at
the angular level.\\
\textbf{Honest scope.}  The termination $L \leq L_\text{max} =
2 l_\text{max}$ is \emph{exact}, not asymptotic.  It is forced by the
Wigner-$3j$ triangle inequality $|l_1 - l_2| \leq L \leq l_1 + l_2$
applied to the angular Gaunt integral at each vertex
(\S~\ref{sec:proj_3j}): for $L > l_1 + l_2$ the $3j$ symbol vanishes
identically, so the multipole sum has no nonzero contribution past
$L_\text{max}$.  There is no truncation parameter to tune, no UV
cutoff to remove, no convergence series to evaluate.  This is
categorically different from the slowly convergent Shibuya--Wulfman
\emph{basis} expansion (Paper~19 \S~III.B), which expands a complete
orbital basis on one center in the Sturmian basis of another center
and does \emph{not} terminate; the multipole expansion of the
$1/|\vec{r}_1 - \vec{r}_2|$ \emph{kernel} terminates trivially while
the Shibuya--Wulfman expansion of the \emph{basis} does not.
Computational note: where the radial integrand at fractional screening
is non-polynomial (the smooth $[Z_A - Z_\text{eff}(r)]/r$ correction
in composed geometries), the algebraic expansion does not apply and
Gauss--Legendre quadrature is used for the radial integral only; the
angular content remains exactly truncated regardless
(Paper~15 \S~V, closing paragraph).\\
\textbf{Structural reading.}  Layer~1 $\to$ Layer~1.  This projection
reorganizes the angular content of the Coulomb kernel \emph{prior to}
any focal-length, mass, or observation projection; it does not move
the chain across the Layer-1/Layer-2 boundary of
\S~\ref{sec:layer2_boundary}.  Categorically, it is the coordinate-side
counterpart of \S~\ref{sec:proj_3j}: the multipole expansion of
$1/|\vec{r}_1 - \vec{r}_2|$ produces the angular vertices at which
Wigner-$3j$ coupling is then applied, so the two projections are
inverse images of each other across the kernel-versus-vertex divide
($P_L(\cos\theta_{12}) \leftrightarrow G(l_1, L, l_2)$).  Operationally
they almost always compose: in the framework's matrix-element pipeline
the multipole expansion is the kernel-side step that produces the
$3j$-symbol calls at each vertex.  Structurally distinct from
\S~\ref{sec:proj_tensor_multipole} (nuclear tensor multipole coupling),
which uses the same Wigner-$3j$ angular machinery but applies it to
nuclear-structure tensor moments (one new $[L]^L$ length scale per
rank); the present projection introduces no physical scale at all and
sits cleanly inside Layer~1.  This is the workhorse projection that
underwrites the entire angular-sparsity story of Paper~22 and the
$O(Q^{2.5})$ Pauli scaling of Paper~14: every claim about ERI density,
selection-rule survival, or composed-architecture sparsity ultimately
rests on the exact $L \leq 2 l_\text{max}$ termination of this
projection.

